\section{Identification of the critical exponent}
\label{sec:critical}

\begin{proof}[Proof of Theorem~\ref{thm:main}]
Apply Proposition~\ref{prop:spin-upper} at \(b=b_c\).
If \(a(b_c)=0\), it gives
\[
 \limsup_n\frac{\log C_{b_c}(n)}{\log n}=-\infty,
\]
whereas Proposition~\ref{prop:critical-lower} gives a liminf of at least
\(-1\). Hence \(a(b_c)>0\), and
Theorem~\ref{thm:coefficient-gap} implies
\[
                         a(b_c)\ge8\pi.
\]
Suppose the inequality were strict.
Proposition~\ref{prop:openness} would give \(a(\widetilde b)>0\)
for all \(\widetilde b\) in a neighborhood of \(b_c\).
Since \(b_c>0\), this neighborhood contains some
\(0<\widetilde b<b_c\).
Proposition~\ref{prop:spin-lower} then gives a polynomial lower
bound for \(C_{\widetilde b}(n)\), and therefore
\(m(\widetilde b)=0\): indeed \(C_{\widetilde b}(n)\le1\), and
its negative logarithm is \(O(\log n)\).
This contradicts \eqref{eq:critical-definition}. We have proved
\[
                              a(b_c)=8\pi.
\]
Propositions~\ref{prop:spin-upper} and~\ref{prop:spin-lower}
now give
\[
 -\frac14
 \le\liminf_n\frac{\log C_{b_c}(n)}{\log n}
 \le\limsup_n\frac{\log C_{b_c}(n)}{\log n}
 \le-\frac14.
\]
Both estimates use the free-box limit first and then all integer
separations, so they establish exactly the limit in the theorem.
The quantified bounds follow from the definition of this limit.
\end{proof}

The conclusion does not assert convergence of
\(n^{1/4}C_{b_c}(n)\). In particular, multiplicative logarithmic
corrections are compatible with the result.
